\section{An exact probability representation of the complete error}
\label{sec:error-law}

From now on fix \(\lambda=0\). The preceding quadrature statements
remain valid at other fixed times, but the primary zero question is
the time-zero one. Define the first omitted Taylor term
\begin{equation}
 L_N(z)=(-1)^N h_N^*\frac{z^{2N}}{(2N)!},\qquad
 \dd\nu_N(u)=\frac{P_N(u)^2}{h_N^*}\,\dd\mu(u).
 \label{eq:leading-error}
\end{equation}
The measure \(\nu_N\) is a probability; it weights the original
source by the squared degree-\(N\) polynomial. The factorial and sign
in \(L_N\) agree with the first missing return path in
\eqref{eq:first-walk}. The following result explains the entire
remainder after this first term.

\begin{theorem}[Complete error law]
\label{thm:error-law}
Let \(U\) have law \(\nu_N\). Independently, choose a vector
\[
 S=(S_0,S_{1,1},S_{1,2},\ldots,S_{N,1},S_{N,2})
\]
uniformly on the simplex of nonnegative coordinates summing to one. Set
\begin{equation}
 X_N=S_0U+\sum_{j=1}^N(S_{j,1}+S_{j,2})x_{j,N},\qquad
 E_N(z)=\mathbb E\e^{\ii zX_N}.
 \label{eq:error-random-variable}
\end{equation}
Then \(X_N\) has a symmetric law, \(E_N\) is even and entire,
\(E_N(0)=1\), and on all of \(\mathbb C\),
\begin{equation}
 \boxed{F(z)-G_N(z)=L_N(z)E_N(z).}
 \label{eq:error-factorization}
\end{equation}
Thus \(F-G_N\) vanishes to order exactly \(2N\) at the origin.
\end{theorem}

The expectation in \eqref{eq:error-random-variable} represents the
error, not a replacement physical source. Each Gauss node appears
twice because the interpolation used in Gaussian quadrature matches
both a value and a derivative. We give the normalization and the
complex-height derivation explicitly; the underlying confluent
interpolation and simplex formulas are classical
\cite[\S1, Eq.~(16); \S7, Definition~37; \S9, Eq.~(52)]{deboor}.

\begin{proof}
Put \(d=2N\). On a simplex with \(d+1\) coordinates, integration
of \(\prod_i s_i^{k_i}\), with \(k=\sum_i k_i\), gives
\(\prod_i k_i!/(d+k)!\). To check this, fix one coordinate \(r\),
rescale the remaining simplex by \(1-r\), and induct on dimension
using \(\int_0^1r^a(1-r)^b\,\dd r=a!b!/(a+b+1)!\).
In particular its coordinate volume is \(1/d!\), so uniform
\emph{probability} has density \(d!\).
For fixed knots \(t_0,\ldots,t_d\), define
\(c_k(t)=\sum_{k_0+\cdots+k_d=k}\prod_i t_i^{k_i}\). The
multinomial formula now gives
\begin{equation}
 \mathbb E\Big(\sum_iS_it_i\Big)^k
   =\frac{k!d!}{(d+k)!}c_k(t),\qquad
 \sum_{k\ge0}c_k(t)v^k=\prod_i(1-vt_i)^{-1}.
 \label{eq:simplex-moments}
\end{equation}

Let \(I_z\) interpolate \(f_z(v)=\e^{\ii zv}\) and its first
derivative at the \(N\) Gauss nodes, in degree at most \(d-1\).
It exists uniquely because homogeneous data would force divisibility
by \(P_N^2\), which has degree \(d\). Adjoin a real source point
\(u\) to the doubled nodes, forming
\(t=(u,x_{1,N},x_{1,N},\ldots,x_{N,N},x_{N,N})\). Interpolation
on this list in degree at most \(d\), with derivative data at
repeated knots, shows
\[
 f_z(u)-I_z(u)=P_N(u)^2D_tf_z,
\]
where \(D_tf_z\) is the leading coefficient of the augmented
interpolant. This equality also holds when \(u\) is a node.

For a monomial \(v^l\), its augmented interpolant is the remainder
modulo \(\prod_i(v-t_i)\). The coefficient of \(v^d\) in that
remainder is zero for \(l<d\), and \(c_{l-d}(t)\) for \(l\ge d\),
by expansion of the quotient at infinity. Since \(D_t\) depends
on finitely many values and derivatives, it can be applied to the
exponential power series term by term. Formula~\eqref{eq:simplex-moments}
then gives
\begin{equation}
 D_tf_z=\sum_{k\ge0}\frac{(\ii z)^{d+k}}{(d+k)!}c_k(t)
       =\frac{(\ii z)^d}{d!}
         \mathbb E_S\e^{\ii z\sum_i S_it_i}.
 \label{eq:confluent-exponential}
\end{equation}
This argument works for complex \(z\); it does not use a
real mean-value remainder for a complex exponential.

Exactness for \(I_z\) and agreement with \(f_z\) at the nodes give
\(F-G_N=\int(f_z-I_z)\,\dd\mu\). Integrating
\eqref{eq:confluent-exponential} against \(P_N^2\dd\mu\) proves
\eqref{eq:error-factorization}. All interchanges are justified on
\(|z|\le R\) by \(\e^{R(|U|+\max_j|x_{j,N}|)}\), which is
integrable at each fixed degree by \eqref{eq:tail}. This also proves
entire dependence. Reflection of \(U\) and simultaneous permutation
of the node pairs sends \(X_N\) to \(-X_N\) without changing its
law. Finally \(E_N(0)=1\) and \(h_N^*>0\) give the stated order.
\end{proof}

\subsection{The size of the error variable}

Put \(M=2N+1\), \(T_j=\operatorname{tr}J_N^j\), and
\(m_j^{[N]}=\mathbb E_{\nu_N}U^j\). We use
\((M)_k=M(M+1)\cdots(M+k-1)\), with \((M)_0=1\).

\begin{proposition}[Second and fourth moments]
\label{prop:error-moments}
\begin{equation}
 \begin{aligned}
 \sigma_N^2:=\mathbb E X_N^2
   &=\frac{2(m_2^{[N]}+T_2)}{(M)_2},\\
 \mathbb E X_N^4
   &=\frac{24m_4^{[N]}+24m_2^{[N]}T_2+12T_2^2+12T_4}{(M)_4},
 \end{aligned}
 \label{eq:error-moments}
\end{equation}
where
\begin{equation}
 \begin{aligned}
 m_2^{[N]}&=\beta_N+\beta_{N+1},\qquad
 T_2=2\sum_{j=1}^{N-1}\beta_j,\\
 m_4^{[N]}&=\beta_N\beta_{N-1}+(\beta_N+\beta_{N+1})^2
                 +\beta_{N+1}\beta_{N+2}.
 \end{aligned}
 \label{eq:source-error-moments}
\end{equation}
At \(N=1\) the term containing \(\beta_0=0\) is absent.
\end{proposition}

\begin{proof}
For the doubled symmetric nodes,
\[
 \prod_j(1-vx_{j,N})^{-2}
   =1+T_2v^2+\tfrac12(T_2^2+T_4)v^4+O(v^6),
\]
by expanding the formal logarithm. At fixed \(U=u\), multiply
by \((1-vu)^{-1}\); the coefficients of degrees two and four
are \(u^2+T_2\) and
\(u^4+u^2T_2+(T_2^2+T_4)/2\). Average these finite polynomials
and apply \eqref{eq:simplex-moments}. No reciprocal series with
unbounded random \(U\) is integrated.
The moments of \(U\) are the squared norms of \(uq_N\) and
\(u^2q_N\). The recurrence gives the coefficients
\(a_N,a_{N+1}\) for the former and
\(a_Na_{N-1},\beta_N+\beta_{N+1},a_{N+1}a_{N+2}\) for the latter.
Orthogonality gives \eqref{eq:source-error-moments}.
\end{proof}

For example, \(N=1\) has one node at zero and
\(X_1=S_0U\), where \(S_0\) has density \(2(1-s)\) on
\([0,1]\). Thus \(\mathbb E X_1^2=m_2^{[1]}/6\) and
\(\mathbb E X_1^4=m_4^{[1]}/15\), checking the endpoint factors.
The envelope already shows
\(\sigma_N^2=O(\log^2(N+1)/N)\). To control complex phase,
we must additionally control an exponential weight; variance alone
does not provide that control.
